% =====================================================================
\chapter{Структура SD-карты}
\label{app:sd}
% =====================================================================

\begin{Code}[numbers=none]
/
├── FIRMWARE/        файлы прошивок (.bin) для загрузчика и модулей
├── LOGS/            логи телеметрии (CSV)
├── MODELS/          модели: model00.yml, model01.yml, ...
├── RADIO/           radio.yml — настройки пульта
├── SCREENSHOTS/     снимки экрана
├── SCRIPTS/
│   ├── FUNCTIONS/   функциональные Lua-скрипты
│   ├── MIXES/       микшерные Lua-скрипты
│   ├── TELEMETRY/   экраны телеметрии
│   ├── TOOLS/       утилиты (elrsV3.lua и др.)
│   └── WIZARD/      мастера создания моделей
├── SOUNDS/
│   ├── en/          английский голосовой пакет
│   │   └── SYSTEM/  системные фразы
│   └── ru/          русский голосовой пакет (если установлен)
├── IMAGES/          картинки моделей
└── EEPROM/          резервные копии (при конвертации старых версий)
\end{Code}

\begin{itemize}
  \item Минимально необходимые папки для работы: \code{RADIO}, \code{MODELS},
        \code{SOUNDS} (для голоса) и \code{SCRIPTS} (для ELRS).
  \item Карта --- microSD, файловая система FAT32. Карт объёмом 8--32\,ГБ более чем достаточно.
  \item Если карта повреждена: отформатируйте её в FAT32, запишите пакет SD через Buddy и
        верните из резервной копии папки \code{MODELS} и \code{RADIO}.
\end{itemize}
